\begin{center}
\LARGE
Generic solutions to polynomial equations in distributive
                            categories
\end{center}

\begin{center}
\Large
Robbie Gates
\end{center}

\noindent
Date:  July 18th (Thursday)\\
Time:  15:25-15:45\\

\begin{center}
\large
Abstract
\end{center}

Following a remark of Lawvere, Blass exhibited a particularly elementary bijection between the
set of binary trees, and seven tuples of binary trees. A "set of binary trees" may be abstracted to
"an object $T$ of a distributive category with a given isomorphism $T^2 + 1 \cong T$".
Particularly elementary in fact means "constructible from the given isomorphism using
distributive category operations". In this context, it is seen that there is no such bijection for pairs,
triples, ..., six-tuples of binary trees. The author has described, for a given polynomial P, the free
distributive category containing an object $X$ and given isomorphism $P(X) \cong X$, and the
isomorphism classes of this category. Since distributive operations correspond to the operations
used to build straightline circuits/programs from simpler circuits/programs, the results of the
author may be interpreted as proving the existence/non-existence of isomorphims constructible by
straight line circuits/programs. The talk will briefly describe the connection between distributive
categories and circuits, the bijection presented by Blass, the techniques used and results acheieved
for the general case, and some speculation on future directions. 

